\Needspace{42\baselineskip}
\section{Bidirectional extension groupoid and operator algebra}
\label{sec:cpe-grupoide-two-way}

Let \(X_F^{\leftrightarrow}\) be the natural extension of the enriched state and
\(\sigma\) its shift homeomorphism. Define the transformation
groupoid
\begin{equation}
 \mathcal G_{\rm TPK}
 :=X_F^{\leftrightarrow}\rtimes_\sigma\mathbb Z,
 \label{eq:cpe-grupoide-tpk}
\end{equation}
with
\[
 s(x,n)=x,
 \quad r(x,n)=\sigma^n x,
 \quad (\sigma^nx,m)(x,n)=(x,m+n),
 \quad (x,n)^{-1}=(\sigma^nx,-n).
\]

\begin{theorem}[Groupoidal Completion of Bidirectional Transport]
\label{thm:cpe-grupoide-tpk}
The groupoid \(\mathcal G_{\rm TPK}\) is locally compact, Hausdorff,
second countable, etale, and amenable. For the invariant and
atomless projective measure \(\mu_{\rm HMT}\), if the shift is ergodic and
essentially free,
\[
 \mathcal M_{\rm TPK}
 :=L^\infty(X_F^{\leftrightarrow},\mu_{\rm HMT})
 \rtimes_\sigma\mathbb Z
\]
is the hyperfinite factor of type \(II_1\).
\end{theorem}

\begin{proof}
The action of a countable discrete group by homeomorphisms on a compact
metrizable space produces an etale, locally compact, Hausdorff, and
second-countable groupoid. Amenability of \(\mathbb Z\) implies amenability
of the groupoid. The projective measure is shift-invariant.
Ergodicity and essential freeness make the crossed product a factor;
nonatomicity excludes the matrix type-I case. The trace induced by
\(\mu_{\rm HMT}\) is finite and faithful, and amenability yields
hyperfiniteness. Therefore, the factor is of type \(II_1\).
\end{proof}

The type \(III\) modular classification belongs to a distinct
quasi-invariant measure, with a declared Radon--Nikodym cocycle. It is not
deduced from the mere presence of KMS dynamics and does not reopen the prior
groupoidal completion.
